\subsection{Byte-Oriented Data Streams}

Binary mode (\texttt{"rb"}, \texttt{"wb"}, \texttt{"r+b"}) expects and returns \texttt{bytes}-like objects. Passing \texttt{str} to \texttt{write()} raises \texttt{TypeError}---encoding must be explicit.

\subsubsection{Binary Mode as Raw Byte Transfer Without Text Decoding}

\begin{verbatim}
path = "binary_mode_demo.bin"
data = b"Binary mode does not decode text. 42.\n"
with open(path, "wb") as f:
    f.write(data)
with open(path, "rb") as f:
    restored = f.read()
print(data == restored, type(restored))
\end{verbatim}

\subsubsection{The \texttt{bytes} Object: Immutable Byte Sequences Distinct from \texttt{str}}

\texttt{bytes} elements are integers \texttt{0}--\texttt{255}; \texttt{str} elements are Unicode codepoints. Indexing \texttt{bytes} yields \texttt{int}; slicing yields \texttt{bytes}.

\subsubsection{The \texttt{bytearray} Object: Mutable Byte Buffers for Incremental Modification}

\texttt{bytearray} supports in-place mutation and growth---useful for incremental assembly before a single binary write.

\begin{verbatim}
buf = bytearray(b"Answer: 41")
buf[-1] = ord("2")
print(buf.decode("utf-8"))
\end{verbatim}
